\section{Closure theorems from the admissibility complex}
\label{sec:closure-theorems}

The preceding two sections introduced the carrier-compatible geometric and transport architecture.
The primitive admissibility projector $\Pprim$ was already placed in the primitive section so
that the proof direction runs from the boundary datum through the derived closed datum to $\Pprim$, and only then to the later
sector theorems. What remains in this section is therefore (i) the upgrade of $\Pprim$ to the
full admissibility projector $\Padm$ once the discrete carrier and determinant data have been
fixed, and (ii) the closure stack that follows from that full projector together with the family,
Higgs, positivity, and reflection-stability inputs collected first.

\begin{definition}[Full admissibility complex]
\label{def:full-admissibility-complex}
Once the discrete carrier and determinant data are fixed on the admissible branch, the Hilbert
space carries a color factor and a determinant involution. Write
\[
\mathcal H=\mathcal H_{\Seam}\otimes\mathcal H_c\otimes\mathcal H_\Theta\otimes\mathcal H_t,
\qquad
\mathcal H_{\mathrm{ext}}
:=
\mathcal H\otimes\Lambda(\eta_\Sigma,\eta_c,\eta_\Theta,\eta_g),
\]
with carrier-side color generator $Z_c$ and determinant-sector involution $Q_\Theta$ satisfying
\[
\iC^2=\mathbf 1,
\qquad
Q_\Theta^2=\mathbf 1,
\qquad
Q_\Theta^\dagger=Q_\Theta.
\]
Define the full structural selector
\[
P_{\mathrm{str}}:=P_{\Seam,+}P_{\mathrm{sing}}P_\Theta,
\qquad
P_{\Seam,+}:=\frac{\mathbf 1+\iC}{2},
\qquad
P_\Theta:=\frac12(\mathbf 1+Q_\Theta),
\]
\[
P_{\mathrm{sing}}:=\frac13(\mathbf 1+Z_c+Z_c^2).
\]
Let $\epsilon_{\mathrm{adm}}>0$ and let $D_{\mathrm{adm}}$ be the admissible sector operator on
$\mathrm{Ran}(P_{\mathrm{str}})$ commuting with $\iC$, $P_{\mathrm{sing}}$, and $Q_\Theta$.
Define the elementary penalties
\[
K_\Sigma:=\sqrt{\lambda_\Sigma}\,\frac{\mathbf 1-\iC}{2},
\qquad
K_c:=\sqrt{\lambda_c}\,\bigl(\mathbf 1-P_{\mathrm{sing}}\bigr),
\]
\[
K_\Theta:=\sqrt{\lambda_\Theta}\,\frac{\mathbf 1-Q_\Theta}{2},
\qquad
K_g:=\sqrt{\lambda_{\mathrm{gap}}}\,\mathbf 1_{[0,\epsilon_{\mathrm{adm}}^2/4)}\!\bigl(P_{\mathrm{str}}D_{\mathrm{adm}}^\dagger D_{\mathrm{adm}}P_{\mathrm{str}}\bigr).
\]
The full admissibility differential is
\[
Q_{\mathrm{adm}}
:=
\eta_\Sigma^\dagger K_\Sigma
+\eta_c^\dagger K_c
+\eta_\Theta^\dagger K_\Theta
+\eta_g^\dagger K_g,
\]
and the full admissibility Laplacian is
\[
\Delta_{\mathrm{adm}}:=\{Q_{\mathrm{adm}},Q_{\mathrm{adm}}^\dagger\}.
\]
\end{definition}

\begin{theorem}[Two-stage upgrade from $\Pprim$ to $\Padm$]
\label{thm:padm-upgrade}
On the closed branch obtained after the carrier theorem and the compact bosonic index closure,
\[
Q_{\mathrm{adm}}^2=0,
\qquad
\Delta_{\mathrm{adm}}=K_\Sigma^2+K_c^2+K_\Theta^2+K_g^2=\Kadm,
\]
and the full Hodge projector
\[
\Padm
:=
\Pi_{\ker\Delta_{\mathrm{adm}}}
=
\operatorname*{s-lim}_{\beta\to\infty}e^{-\beta\Kadm}
\]
admits the two-stage factorization
\[
\Padm
=
\Pprim\,P_{\mathrm{sing}}\,P_\Theta
=
P_{\mathrm{str}}\,P_{\mathrm{gap}},
\qquad
P_{\mathrm{gap}}
:=
\mathbf 1_{[\epsilon_{\mathrm{adm}}^2/4,\infty)}\!\bigl(P_{\mathrm{str}}D_{\mathrm{adm}}^\dagger D_{\mathrm{adm}}P_{\mathrm{str}}\bigr).
\]
The first factor $\Pprim$ depends only on the primitive datum
$(\iC,\APSdom,D_{\mathrm{rel}})$; the second factor $P_{\mathrm{sing}}P_\Theta$ depends only on
the carrier and determinant data made available by the carrier and bosonic index theorems.
\end{theorem}

\begin{proof}[Proof sketch]
The auxiliary generators $\eta_\Sigma^\dagger,\eta_c^\dagger,\eta_\Theta^\dagger,\eta_g^\dagger$
anticommute, so $Q_{\mathrm{adm}}^2$ reduces to the commutators of $K_\Sigma,K_c,K_\Theta,K_g$.
The four commuting sectors $\mathcal H_\Seam,\mathcal H_c,\mathcal H_\Theta,\mathcal H_t$ are
acted on independently by these penalties, hence the mixed terms vanish and $\Delta_{\mathrm{adm}}=\Kadm$. The strong-coupling limit therefore projects onto the harmonic
admissible sector. Because $K_\Sigma$ and $K_g$ already define $\Pprim$ on
$\mathcal H_{\mathrm{prim}}$ and the additional penalties $K_c,K_\Theta$ commute with both, the
full projector factorizes through $\Pprim$ multiplicatively, giving the two-stage form.
\end{proof}

\begin{corollary}[Factorized selector as a derived form]
\label{cor:factorized-admissibility-selector}
On the admissible branch the familiar factorized selector
\[
\Padm=P_{\Seam,+}P_{\mathrm{sing}}P_\Theta\,P_{\mathrm{gap}}
\]
is recovered as a corollary of \Cref{thm:padm-upgrade}; nothing in the primitive section depends
on it.
\end{corollary}

The two-stage selector geometry is summarized in \Cref{fig:admissibility-stack}.

\begin{figure}[t]
\centering
\scriptsize
\begin{tikzpicture}[
    >=Latex,
    data/.style={draw=black!60, fill=gray!8, rounded corners, text width=2.55cm, minimum height=0.9cm, align=center},
    prim/.style={draw=blue!60!black, fill=blue!6, rounded corners, text width=2.55cm, minimum height=0.95cm, align=center},
    close/.style={draw=orange!75!black, fill=orange!8, rounded corners, text width=2.55cm, minimum height=0.95cm, align=center},
    result/.style={draw=green!50!black, fill=green!8, rounded corners, text width=3.55cm, minimum height=1.0cm, align=center},
    note/.style={font=\footnotesize\itshape, text=black!65, align=center}
]
\node[data] (seamdata) at (-3.3,1.95) {primitive datum\\$\iC,\ \APSdom,\ D_{\mathrm{rel}}$};
\node[data] (carrierdata) at (3.3,1.95) {carrier \& determinant closure\\$X,\ Z_c,\ Q_\Theta$};

\node[prim] (pseam) at (-4.55,0.15) {$P_{\Seam,+}$\\seam-even sector};
\node[prim] (pgap) at (-2.0,0.15) {$P_{\mathrm{gap}}$\\admissible gap};
\node[prim] (pprim) at (-3.3,-1.7) {$\Pprim$\\primitive selector};

\node[close] (psing) at (2.0,0.15) {$P_{\mathrm{sing}}$\\color singlets};
\node[close] (ptheta) at (4.55,0.15) {$P_\Theta$\\determinant sector};
\node[close] (pupgrade) at (3.3,-1.7) {$P_{\mathrm{sing}}P_\Theta$\\closure upgrade};

\node[result] (padm) at (0,-3.7) {$\Padm=\Pprim P_{\mathrm{sing}}P_\Theta$\\full admissibility projector};

\draw[->, thick] (seamdata) -- (pseam);
\draw[->, thick] (seamdata) -- (pgap);
\draw[->, thick] (carrierdata) -- (psing);
\draw[->, thick] (carrierdata) -- (ptheta);
\draw[->, thick] (pseam) -- (pprim);
\draw[->, thick] (pgap) -- (pprim);
\draw[->, thick] (psing) -- (pupgrade);
\draw[->, thick] (ptheta) -- (pupgrade);
\draw[->, thick] (pprim) -- (padm);
\draw[->, thick] (pupgrade) -- (padm);
\node[note] at (0,-5.0) {Primitive data determine $\Pprim$; carrier and determinant closure upgrade it to $\Padm$.};
\end{tikzpicture}
\caption{Two-stage upgrade from the primitive selector to the full admissibility projector. The
left lane is fixed by the primitive datum, while the right lane enters only after carrier and
determinant closure.}
\label{fig:admissibility-stack}
\end{figure}

\begin{theorem}[Calder\'on upgrade of the APS boundary domain]
\label{thm:apsdom-upgrade}
Let $\APSdomprov$ be the provisional APS realization of the geometric anchor used in
\Cref{thm:well-posed-primitive-dynamics}, and let $D_b$ be the canonical doubling of
$D_{\mathrm{rel}}$ across $\Seam$. Let $C(D_b)$ denote the Calder\'on projector of $D_b$ on the
seam collar. Then there is a canonical boundary-elliptic upgrade
\[
\APSdomprov
\;\xrightarrow{\;C(D_b)\;}\;
\APSdom
\]
such that:
\begin{enumerate}[label=(\roman*), leftmargin=1.5em, topsep=3pt, itemsep=2pt]
    \item $\APSdom$ lies in the same admissible elliptic class as $\APSdomprov$ and preserves the
    self-adjoint Fredholm realization of the geometric anchor;
    \item the relative graph domain of $D_{\mathrm{rel}}$ is unchanged, so all primitive results
    that depended only on $\APSdomprov$ continue to hold with $\APSdom$;
    \item the family-pairing data $\APSfam$ are induced canonically on the doubled admissible
    branch and are not a separate primitive datum.
\end{enumerate}
\end{theorem}

\begin{proof}[Proof sketch]
The Calder\'on projector $C(D_b)$ is canonically associated with the doubled operator on the
seam collar and produces a boundary-elliptic projection compatible with the provisional APS
choice. Both $\APSdomprov$ and $\APSdom$ are admissible elliptic boundary conditions for the
same Dirac-type operator and therefore define equivalent self-adjoint Fredholm realizations of
the geometric anchor (boundary-elliptic equivalence preserves the Sobolev graph domain).
Consequently every statement made on the primitive scaffold under $\APSdomprov$ persists on the
closed branch under $\APSdom$. The induction of $\APSfam$ from canonical doubling is a standard
consequence of the cobordism structure of the doubled APS pair.
\end{proof}

\begin{proposition}[Master closure roadmap from canonical doubling]
\label{thm:master-closure}
\label{prop:master-closure-roadmap}
\label{ass:remaining-closure-inputs}
This proposition is a roadmap statement: it lists the closure conclusions that the individual
theorems of this section establish from $\Tbound$, the derived closed datum, the carrier theorem, and the
canonical doubling. It is not used as an input to those proofs and is recovered as a master
corollary at the end of the section. Let $X_b$ be the canonical double across $\Seam$ and write
\[
D_b=
\begin{pmatrix}
0&D_{\mathrm{rel}}^-\\
D_{\mathrm{rel}}^+&0
\end{pmatrix},
\qquad
Q_{\mathrm{tot}}:=Q_{\mathrm{adm}}+Q_{\mathrm{geo}},
\qquad
\Delta_{\mathrm{tot}}:=\{Q_{\mathrm{tot}},Q_{\mathrm{tot}}^\dagger\}.
\]
Then on the closed branch all of the following are derived from the same datum (each item is the
output of a downstream theorem in this section, not an extra hypothesis here):
\begin{enumerate}[label=(\arabic*), leftmargin=1.5em, topsep=3pt, itemsep=2pt]
    \item the Calder\'on projector of $D_b$ canonically upgrades the provisional APS choice
    $\APSdomprov$ to $\APSdom$ (cf.\ \Cref{thm:apsdom-upgrade});
    \item the family pairing data are induced canonically, so no separate $\APSfam$ datum remains;
    \item the upgrade theorem \Cref{thm:padm-upgrade} produces
    \[
    \Padm:=\Pi_{\ker \Delta_{\mathrm{adm}}}
    =
    \Pprim\,P_{\mathrm{sing}}\,P_\Theta,
    \]
    and the doubled total closure package preserves $\mathrm{Ran}(\Padm)$;
    \item
    \[
    \chi(X_f^\circ)=-2,
    \qquad
    F\simeq H^1(X_f^\circ,\mathbb C)\simeq \mathbb C^3,
    \qquad
    N_{\mathrm{fam}}=3;
    \]
    \item the harmonic determinant line is parallel and therefore
    \[
    \mathrm{Hol}(\nabla_F)\subset SU(3)_F;
    \]
    \item admissible bosonic, fermionic, and geometric reflection positivity hold on $\mathrm{Ran}(\Padm)$;
    \item the physical geometric sector takes the form
    \[
    \mathcal H_{\mathrm{phys}}
    \cong
    \mathcal H_2\otimes\mathcal H_{\mathrm{sc}}\otimes\mathcal H_{\mathrm{matt}},
    \]
    with the scalaron carried by the positive $(\chi_{\mathrm{geo}},R^2)$ block;
    \item $\Padm\Hrel\Padm$ is lower bounded and coercive.
\end{enumerate}
\end{proposition}

\begin{definition}[Relative generating functional]
\label{def:relative-generating-functional}
The relative Euclidean generating functional is normalized as
\[
Z_{\mathrm{rel}}[J,\eta,\bar\eta]
:=
\frac{Z_{\mathrm{adm}}[J,\eta,\bar\eta]}{Z_{\mathrm{ref}}[0,0,0]}.
\]
The reference sector is therefore a fixed normalizing denominator rather than a second fluctuating measure that would mix with admissible positivity.
\end{definition}

\begin{lemma}[Admissible reflection positivity]
\label{lem:admissible-reflection-positivity}
Assume
\[
[\Theta,P_\Theta]=0,
\qquad
\Theta\Padm=\Padm\Theta,
\]
the admissible bosonic kernel satisfies the Bernstein-positivity remark above and is
reflection-stable on the admissible sector, and the
admissible fermionic CAR kernel is reflection positive on that same sector. Then every
admissible observable $O$ supported in $\tau>0$ satisfies
\[
\langle \Theta O,O\rangle_{\mathrm{rel}}
=
\frac{1}{Z_{\mathrm{ref}}[0,0,0]}
\langle \Theta O,O\rangle_{\mathrm{adm}}
\ge 0.
\]
\end{lemma}

\begin{proof}
By definition,
\[
\langle \Theta O,O\rangle_{\mathrm{rel}}
=
\frac{1}{Z_{\mathrm{ref}}[0,0,0]}\,
\langle \Theta O,O\rangle_{\mathrm{adm}}.
\]
The denominator is a fixed positive normalizing constant, so the sign is determined entirely by the
admissible numerator.

Split $O$ into its bosonic and fermionic parts on the retained admissible sector. For the bosonic
part, the Bernstein representation of the admissible kernel writes the Euclidean covariance as a
positive superposition of reflection-stable heat kernels. Reflection across $\tau=0$ therefore sends
every positive-time bosonic observable to a positive quadratic form, so the bosonic contribution to
\[
\langle \Theta O,O\rangle_{\mathrm{adm}}
\]
is nonnegative.

For the fermionic part, the CAR reflection-positivity hypothesis on the admissible subspace implies
that the reflected two-point kernel is positive in the standard fermionic sense; equivalently, the
corresponding Pfaffian / determinant form on positive-time test spinors is nonnegative. Because the
bosonic and fermionic admissible factors commute inside the normalized expectation, their product
remains nonnegative. Dividing by the fixed reference normalization preserves that sign. Hence
\[
\langle \Theta O,O\rangle_{\mathrm{rel}}\ge 0.
\]
\end{proof}

\begin{theorem}[Operative admissibility selector]
Under \Cref{def:full-admissibility-complex,thm:padm-upgrade,prop:master-closure-roadmap}, the
operator $\Padm$ is idempotent on the retained branch and realizes the admissibility predicate
on all sectors used by the closed output statements:
\[
\mathsf{Adm}(X)=1
\quad\Longleftrightarrow\quad
\Padm X = X.
\]
\end{theorem}

\begin{proof}[Proof sketch]
By \Cref{thm:padm-upgrade}, $\Padm=\Pprim P_{\mathrm{sing}}P_\Theta$ is the orthogonal projector
onto the harmonic admissible sector of the full complex and is therefore idempotent by
construction. \Cref{cor:factorized-admissibility-selector} recovers the explicit factorized
selector. The predicate and operator languages agree once the positive transport sector is fixed.
\end{proof}

\begin{theorem}[Topological family closure on the admissible branch]
By \Cref{thm:master-closure,thm:topological-family-closure,cor:boundary-bulk-occupancy},
\[
\chi(X_f^\circ)=-2,
\qquad
F\cong\mathbb C^3,
\qquad
N_{\mathrm{fam}}=3.
\]
Consequently
\[
\Oadm=N_{\mathrm{fam}}\dim S^+=3\times 16=48,
\qquad
\deltatop=\Oadm\,\cthree^4.
\]
\end{theorem}

\begin{proof}[Proof sketch]
This theorem is now a direct consequence of the master closure theorem from Section~6. The
topological family closure and the orbifold corner count have already been compressed there
into one statement, so the present layer merely records the corresponding admissibility
consequence.
\end{proof}

\begin{theorem}[Bosonic index and local spectral-scale closure]
On the admissible branch,
\[
\Ind_{\mathrm{rel}}^{\Padm}(\mathcal B_{E_2}^+)=1,
\qquad
\Ind_{\mathrm{rel}}^{\Padm}(\mathcal B_{E_3}^+)=0,
\]
so the unique light seam-even bosonic zero mode is
\[
\Phi\in\Gamma(E_2)\cong(1,2)_{1/2},
\qquad
\taudbl^{\ast}\Phi=+\Phi.
\]
Moreover the same closure layer yields
\[
\begin{aligned}
\frac{\Mpl^2}{\lambda_\Sigma^2}&=\frac{\rho_\star}{2\pi^2},
&
G_N\lambda_\Sigma^2&=\frac{\pi}{4\rho_\star},
\\
\frac{\vgeo}{\Mpl}&=\gcar\,\beta_{\mathrm{rad}}^2
\exp\!\left[-\frac{\alpha^{-1}(0)+\delta_{\mathrm{ph}}}{5}\right],
\\
G_N\vgeo^2&=\frac{1}{8\pi}\gcar^2\beta_{\mathrm{rad}}^4
\exp\!\left[-\frac{2(\alpha^{-1}(0)+\delta_{\mathrm{ph}})}{5}\right].
\end{aligned}
\]
Here $\vgeo$ is the seam-even UV source scale. The physical electroweak benchmark is the residue
matched quantity $v_{\mathrm{phys}}$ of \Cref{thm:ew-geometric-to-physical-matching}.
\end{theorem}

\begin{proof}[Proof sketch]
The bosonic sector is treated as a genuine Callias / Dirac / Dolbeault-type index problem rather
than a scalar counting exercise. Relative heat-kernel closure together with
\Cref{prop:canonical-bosonic-normalization,lem:two-sheet-zero-mode-normalization,thm:absolute-spectral-planck-closure}
produces the boundary-normalized Einstein branch, the canonical Higgs normalization, and the
seam-even zero-mode formula for $\vgeo$.
\end{proof}
